\section*{ID9801: Relation: (Log n)\textasciicircum{}π}
\paragraph{Metadata.}
PiID: 0201410206723 | RTEID: RTE:P38059.4596:T6.1616 | UUID: fedbddd1-598a-5eb7-bfe2-61510ec278b7

\paragraph{Canonical Statement.}
es, and overheads increase sub-polynomially. Thus, adding more qubits actually accelerates the system’s effective power (contrary to some approaches where adding more qubits hits diminishing returns due to noise). The metamaterial approach means we don’t have to route microwave lines or laser beams to each of 153600 basis states – they are inherently connected via the material’s modes. This avoids a notorious bottleneck (control signal fan-out) that many quantum computers have. To be concrete: if a problem requires doubling the number of logical qubits, our overhead might go up by a factor of \$(\textbackslash{}log n)\textasciicircum{}\textbackslash{}pi\$. If \$n\$ from 100 to 200, log2(100)\textasciitilde{}6.64, log2(200)\textasciitilde{}7.64, ratio \$(7.64/6.64)\textasciicircum{}\textbackslash{}pi \textbackslash{}approx (1.15)\textasciicircum{}\{3.14\} \textbackslash{}approx 1.17\$. So 17\% more physical qubits needed for double the logical qubits – an extremely favorable scaling. Similarly, solving a problem twice as large might only increase the time by a factor of \$2\textasciicircum{}\{1/\textbackslash{}pi\}\textbackslash{}approx2\textasciicircum{}\{0.318\}=1.25\$, i.e., 25\% more time for double size. This is almost an exponential speedup pattern, suggesting an eventual breakeven point beyond which classical computing is ho

\paragraph{Canonical Equation.}
\[
(\log n)^\pi
\]

\paragraph{Dictionary.}
Relation: (Log n)\textasciicircum{}π — (log n)\textasciicircum{}π

\paragraph{Encyclopedia.}
Relation: (Log n)\textasciicircum{}π is a symbolic expression, written (log n)\textasciicircum{}π. It introduces a symbol or relation used elsewhere in the framework. It is built from π, n. Within the Trawin Zero Point registry it serves as one element of the framework's mathematical narrative.
